\section{Proof of the main theorem}\label{sec:main-proof}
This appendix proves Theorem~\ref{thm:main}, Theorem~\ref{thm:sign}, and Corollary~\ref{cor:balanced}.
Section~\ref{sec:proof-sketch} first gives the proof sketch; the subsequent subsections give the complete
chain of proof, and the expansion of the key estimates is collected in Appendix~\ref{app:technical}.

The proof proceeds in five steps. The first step analyzes the initial degree distribution and obtains the
joint Gaussian limit of the energy and the influence function (Section~\ref{sec:mp-init}). The second step
tracks the full degree trajectory along the peeling process, constructs the precritical score, and decomposes
it into an initial-state part and the dynamic innovation (Sections~\ref{sec:mp-traj} and \ref{sec:mp-stop}).
The third step proves that the sign of the score determines the actual success or failure
(Section~\ref{sec:mp-global}). The fourth step transfers the result from the configuration model to the
Plain model (Section~\ref{sec:mp-plain}). The fifth step handles general sign composition
(Section~\ref{sec:signproof}). The two subsections after the proof sketch first give the critical constants
and then fill in the details of a few short arguments in the main text.

\subsection{Proof sketch}\label{sec:proof-sketch}
The proof of Theorem~\ref{thm:main} first establishes the joint limit of $(X_{\rm one},R_M)$ in the
all-equal-sign case, and then handles general sign composition via Theorem~\ref{thm:sign}. The
all-equal-sign case proceeds in four steps; the complete argument appears in the subsequent subsections, and
the key estimates appear in Appendix~\ref{app:technical}.

\paragraph{Initial state.} We first work in the configuration model: the $k$ endpoints of each edge are
drawn independently and uniformly, with repetition allowed. The measurement depends only on the initial
degrees, so it can be written as a functional $G_M(e_M)$ of the degrees, where $e_M(j)=(j-\mu_M)^2$. The
initial-state part of the score is another functional $G_M(h_c)$, with influence function
$h_c(j)=j(1-a_c)^{j-1}$ depending on the complete profile of the initial degrees. The total number of
endpoints is fixed, so the joint Gaussian limit of the two must subtract the correction term in
\eqref{eq:bridge}.

\paragraph{Peeling trajectory and the score.} Along the peeling process, the degree histogram follows the
fluid equation; at the critical parameters the degree-one density has a non-degenerate quadratic minimum,
which is the bottleneck of the peeling. The score $R_M$ reads the histogram at the entry time before the
bottleneck and extrapolates to the bottleneck with backward weights (\eqref{eq:score}). A Doob decomposition
of the backward functional writes $R_M$ as three terms: the deterministic term $\beta_kr$ from the window
location, the initial-state part $G_M(h_c)/\mu_c$, and the dynamic innovation $W$ (\eqref{eq:decomp}). The
increments of the innovation are $O(M^{-1/2})$, and its quadratic variation tends to the constant
$\tau_k^2$; stable convergence (\eqref{eq:stable}) shows that it is independent of the initial-state limit.
Hence the correlation between the measurement and the score comes entirely from the initial state, equal to
$\Gamma(e,h_c)/(\sqrt2\mu_c^2\sigma_k)$. Substituting $\Gamma(e,h_c)=\mu_ce^{-x_c}x_c(x_c-2)$ gives
\eqref{eq:eta}, and the sign of $\eta_k$ is exactly the sign of $x_c-2$. The influence function is
determined by the complete degree profile; tracking only the two lowest-order statistics of the initial
state gives the wrong value $\eta_3\approx-0.398$.

\paragraph{The sign of the score determines success.} Within the short window near the bottleneck the local
noise freezes, and the number of degree-one vertices is approximately $x_c\sqrt M R_M$ plus the fluid
difference (\eqref{eq:freeze}). When $R_M$ is clearly negative, the number of degree-one vertices drops to
zero before the central time, so peeling fails; when $R_M$ is clearly positive, the process crosses the
bottleneck, after which the positive gap of the fluid and a first-moment bound on the small core rule out
failure. The Gaussian limit has no atoms, which gives \eqref{eq:transfer}.

\paragraph{Transfer to the Plain model.} The Plain model is the configuration model conditioned on the
event of having no bad edges, whose probability tends to $e^{-\Lambda_k}>0$. Planting finitely many bad
edges does not change the limit of $(X_{\rm one},R_M)$; expanding with mixed factorial moments and
inclusion-exclusion keeps the joint Gaussian law unchanged after conditioning.

\paragraph{General sign composition.} For statistics invariant under permutation of the element labels, one
can first generate the unsigned graph $\mathcal G$ and then distribute a given sign composition uniformly
over it. Theorem~\ref{thm:sign} decomposes the measurement into $\gamma X_{\rm one}$ and a term $V$
asymptotically independent of the graph. The correlation of the former with the score is $\eta_k$, and
multiplying by $\gamma\to\rho^2$ gives the correlation coefficient $\eta_k\rho^2$. The proof of
Theorem~\ref{thm:sign} itself is in Section~\ref{sec:signproof}: first approximate with iid signs, then flip
a whole batch from the surplus sign class to correct the total sign count exactly, and finally obtain the
normal limit from the self-averaging of the overlap graph and cumulant estimates on the dependency graph.

\subsection{Critical constants}\label{app:constants}
Fix $k\ge3$, and let $x_c$ be the unique positive root of the equation $e^{x}-1=(k-1)x$\cite{molloy,dn,jl07}. Let
\begin{equation*}
q_c=1-e^{-x_c},\quad a_c=q_c^{k-1},\quad
\mu_c=x_c/a_c,\quad\alpha_c=\mu_c/k.
\end{equation*}
$\mu_c$ is the critical average degree and $\alpha_c$ the critical load. In Theorem~\ref{thm:main} the slope
of the score against the window coordinate and the limiting variance of the score are
\begin{equation*}
\beta_k=-ka_ce^{-x_c},\qquad
\sigma_k^2=e^{-x_c}q_c\left(1-\frac{k a_c}{k-1}\right).
\end{equation*}
$\sigma_k^2$ is composed of the initial-state part $\sigma_0^2$ and the dynamic part $\tau_k^2$; see
\eqref{eq:tau} and the simplification that follows it. For $k=3,4,5,6$ the values of $\tau_k^2$ are
$0.031017361$, $0.016877142$, $0.011169746$, $0.008107775$; the numerical values of $\sigma_k^2$ and
$\eta_k$ appear in Table~\ref{tab:constants}.

\paragraph{Finite-size shift constants.}\label{app:ck}
The shift constants of Section~\ref{sec:rules} are obtained by converting the critical threshold shift of
Dembo--Montanari\cite{dm} to the scale of this paper. Write $w_k=\sigma_k/|\beta_k|$,
$\widetilde G=(k-2)/(k-1)$, $\widetilde F=[(k-1)x_c-(k-2)]/(kq_c^k)$. Then
\begin{equation*}
C_k=\frac{1}{w_k}\cdot\frac{0.996193\,\widetilde G^{2/3}\widetilde F^{-1/3}\alpha_c^{4/3}}{x_c^2e^{-x_c}},
\end{equation*}
where the constant $0.996193$ is the numerical value of the Airy integral in DM's eq. (2.16).

\subsection{Exact normalization and local uniformity}\label{app:normalization}
The normalization in \eqref{eq:measurement} of Section~\ref{sec:model} is exact, for the following reason.
Let $H_{ij}=|S_i\cap S_j|$. Then
\begin{equation*}
T=kd(1-k/M)+2\sum_{i<j}s_i s_j(H_{ij}-k^2/M).
\end{equation*}
The hypergeometric distribution gives
\begin{equation*}
\E H_{ij}=k^2/M,\qquad
\Var H_{ij}=\frac{k^2(M-k)^2}{M^2(M-1)}.
\end{equation*}
Different unordered pairs of elements have uncorrelated overlaps. Even when two pairs share an element, once
the mapping of that element is given, the conditional mean of each pair is still the constant $k^2/M$.
Summing these quantities gives the normalization in \eqref{eq:measurement}.

\paragraph{Local uniformity.} The window assumption \eqref{eq:window} allows an $o(M^{-1/2})$ slack, and the
way $\rho_M\to\rho$ is unrestricted, so any sequence with $r_M\to r$ satisfies the conditions of the theorem.
A subsequence argument gives: the joint limit in Theorem~\ref{thm:main} and the limit of the conditional CDF
hold uniformly over $(r,\rho)$ on any compact subset of $\mathbb R\times[-1,1]$. This removes the need to
prove a separate uniform CLT.

\paragraph{Convergence of conditional moments.} Convergence in distribution alone does not imply convergence
of moments: a small part of the probability mass can escape to infinity, vanish in the limit, yet still bias
the mean. Section~\ref{sec:mean} therefore needs $X^2$ to be uniformly integrable under failure. This has two
layers. The first is unconditional. Since $X\Rightarrow N(0,1)$ and $\E X^2=1$ holds exactly for every $M$,
the second moment equals the second moment of the limit, so $X^2$ is uniformly integrable. The second is
conditional. For any $K>0$,
\begin{equation*}
\E\bigl[X^2;\,X^2>K\bigm|F\bigr]\le\frac{\E\bigl[X^2;\,X^2>K\bigr]}{\Pr(F)}.
\end{equation*}
The numerator tends to zero uniformly in $M$ as $K\to\infty$, and the denominator tends to $\Phi(z_r)>0$, so
$X^2$ is also uniformly integrable under failure. The critical window requires the limiting failure
probability to be positive, and it matters here: if the denominator tended to zero, conditioning could blow up
the tail without control. Combined with local uniformity, the limits of the first two conditional moments
hold uniformly over $(r,\rho)$ on compact sets.

\subsection{Initial joint Gaussian limit}\label{sec:mp-init}
First relax the model: the $k$ endpoints of each edge are drawn independently and uniformly, with repetition
allowed. This is the configuration model. Write $D_v$ for the initial degree of vertex $v$ and
$\mu_M=kd/M$. For $e_M(j)=(j-\mu_M)^2$ and $h_c(j)=j(1-a_c)^{j-1}$ (with $h_c(0)=0$), define
\begin{equation*}
G_M(f)=\frac1{\sqrt M}\left[\sum_v f(D_v)-M\E f(\Pois(\mu_M))\right].
\end{equation*}
Conditioned on the total number of endpoints being fixed, the occupancy CLT gives the joint Gaussian limit of
$(G_M(e_M),G_M(h_c))$ with covariance
\begin{equation}
\Gamma(f,g)=\Cov(f(D),g(D))-
\frac{\Cov(f(D),D)\Cov(g(D),D)}{\mu_c},\quad D\sim\Pois(\mu_c).\label{eq:bridge}
\end{equation}
The second term is the correction coming from the constraint that the total number of endpoints is fixed.

The concrete approach is de-Poissonization. On the same independent endpoint sequence, replace the fixed
length $L=kd$ by an independent random length $N\sim\Pois(L)$. Write $S_f(l)=\sum_v f(D_v(l))$. Then
\begin{equation*}
S_f(N)-S_f(L)=a_f(N-L)+o_p(\sqrt M),\qquad
a_f=\frac{\Cov(f(D),D)}{\mu_c}.
\end{equation*}
For $e_M$, each newly added endpoint changes the sum on average by $2(l/M-\mu_M)+1$; for $h_c$, both the
first and second discrete differences are bounded. Over the range $|l-L|\le\sqrt M\log M$ the accumulated
drift error is $O_p(\log^2M)$ and the added martingale fluctuation is $o_p(\sqrt M)$. The ordinary Poisson
vector CLT plus this correction gives \eqref{eq:bridge}. The one-sided measurement satisfies
\begin{equation*}
X_{\rm one}=\frac{G_M(e_M)}{\sqrt2\mu_c}+o_p(1).
\end{equation*}
The measurement in the configuration model differs from the center used here, but this difference vanishes on
the $\sqrt M$ scale.

\subsection{Degree trajectory and the precritical score}\label{sec:mp-traj}
Let $t$ be the number of deleted edges divided by $M$, and write
\begin{equation*}
m_M(t)=\mu_M-kt,\quad a_M(t)=\left(\frac{m_M(t)}{\mu_M}\right)^{(k-1)/k},\quad
q_M(t)=a_M(t)^{1/(k-1)}.
\end{equation*}
Write $U_j$ for the number of vertices of current degree $j$ and $Q=U_1$ for the number of degree-one
vertices. While peeling has not stopped, the coordinates of degree $j\ge2$ satisfy the following fluid
equation; starting from the Poisson initial state, its solution is
\begin{equation*}
\dot u_j=\frac{k-1}{m_M(t)}[(j+1)u_{j+1}-ju_j],\qquad
u_j(t)=e^{-\mu_Ma_M(t)}\frac{[\mu_Ma_M(t)]^j}{j!},\quad j\ge2.
\end{equation*}
The corresponding degree-one density is
\begin{equation*}
\ell_M(t)=m_M(t)-\mu_Ma_M(t)[1-e^{-\mu_Ma_M(t)}].
\end{equation*}
At the critical parameters, the interior zero of $\ell_M$ is a non-degenerate quadratic minimum. This is the
bottleneck of the peeling.

Now define the backward weights of the heavy stubs at the final point:
\begin{equation*}
f_j(b)=jb[1-(1-b)^{j-1}],\quad f_0=f_1=0.
\end{equation*}
When the thinning parameter $a_c$ of the final point is fixed, this weight satisfies the backward equation.
The projection of the initial fluctuation onto it is $-\sum_j f_j(a_c)G_j=a_c\sum_jh_c(j)G_j$, where we used
the total constraint $\sum_jjG_j=0$. We adopt the convention that failure corresponds to a low score; under
this convention the direction of the initial-state influence is the \emph{positive} $h_c/\mu_c$.

Take
\begin{equation*}
t_{*,M}=\frac{\mu_M}{k}(1-q_c^k),\quad
\varepsilon_M=M^{-1/8},\quad t_-=t_{*,M}-\varepsilon_M,\quad
n_-=\lfloor Mt_-\rfloor.
\end{equation*}
If the process reaches $n_-$, let $b_-=a_c/a_M(t_-)$ and define
\begin{equation}
R_M=\frac{M m_M(t_{*,M})-\sum_{j\ge2}f_j(b_-)U_j(n_-)}{x_c\sqrt M}.\label{eq:score}
\end{equation}
If the process does not reach $n_-$, set $R_M=0$ for now; we will show later that this case has probability
$o(1)$. $R_M$ is an extrapolated score, computed directly from the degree histogram at the entry time $n_-$,
and it can be negative.

\subsection{Concentration before the bottleneck and dynamic innovation}\label{sec:mp-stop}
\begin{lemma}[Uniform concentration before stopping]
Let $\theta_M$ be the first step at which the number of degree-one vertices is zero. For fixed $\delta>0$,
take $\vartheta_M=\sqrt M(\log M)^6$. Then on $n\le\theta_M\wedge(d-\delta M)$, the uniform errors of
$Q_n-M\ell_M(n/M)$ and of $U_j(n)-M u_j(n/M)$ ($2\le j\le\lceil\log M\rceil$) are all $O_p(\vartheta_M)$.
\end{lemma}
The argument uses the backward binomial weights and $f_j$, and stops the process at $\theta_M$. The maximum
initial degree is at most $L_M=\lceil\log M\rceil$ with high probability. As long as the number of remaining
stubs is $\Theta(M)$, the backward equation cancels the main drift, and the per-step remainder is controlled
by $C_\delta L_M^4/M$, with the projected martingale difference uniformly bounded. Combining three tools
gives the conclusion: the Azuma concentration inequality, the bounded-difference concentration of the initial
state, and a union bound over $O(ML_M)$ target times and degrees. The fluid gap at the entry is
$M\varepsilon_M^2=M^{3/4}$, much larger than $\vartheta_M$, so the process reaches $n_-$ with probability
$1-o(1)$. The detailed estimates are in Appendix~\ref{app:technical}.

A Doob decomposition of the backward functional gives
\begin{equation}
R_M=\beta_k r+\frac{G_M(h_c)}{\mu_c}+W+o_p(1).\label{eq:decomp}
\end{equation}
On the right are, in order, the deterministic term from the window location, the initial-state part, and the
dynamic innovation $W$. The normalized increments of $W$ are $O(M^{-1/2})$, and its predictable quadratic
variation tends to the deterministic constant $\tau_k^2$. Let $\Psi_M$ be measurable with respect to the
initial-state information and uniformly bounded. The stable decomposition we need is
\begin{equation}
\E[\Psi_Me^{iuW}]-e^{-u^2\tau_k^2/2}\E[\Psi_M]\longrightarrow0.\label{eq:stable}
\end{equation}
The proof applies a one-step Taylor expansion to the compensated exponential, with the sum of the remainders
$O(M^{-1/2})$; together with the uniform boundedness of the quadratic variation and its convergence to the
deterministic value, this gives \eqref{eq:stable}. Taking $\Psi_M$ to be the characteristic function of the
initial-state vector yields simultaneously the joint Gaussian limit and the independence between the
innovation and the initial-state limit. The detailed estimates are in Appendix~\ref{app:technical}.

Next we compute the dynamic variance. On the critical fluid trajectory let $b=a_c/a(t)$ and
$g_j(b)=f_j(b)-f_{j-1}(b)$, and set $g_1=0$. A random partner hits a vertex of degree $j\ge2$ with
probability $p_j(t)=j u_j(t)/m(t)$; the remaining probability mass goes into $p_1$. Each step deletes $k-1$
fixed partners, so
\begin{equation}
\tau_k^2=\frac{k-1}{x_c^2}\int_0^{t_c}
\left[\sum_jp_j(t)g_j(b)^2-\left(\sum_jp_j(t)g_j(b)\right)^2\right]dt.\label{eq:tau}
\end{equation}
Let $s_c=1-a_c$. The initial-state contribution $\sigma_0^2$ and the total variance simplify to
\begin{align*}
\sigma_0^2&=(1/\mu_c+s_c^2)e^{\mu_c(s_c^2-1)}
-e^{-2x_c}[1+(1-x_c)^2/\mu_c],\\
\sigma_k^2&=\sigma_0^2+\tau_k^2=e^{-x_c}q_c\left(1-\frac{k a_c}{k-1}\right).
\end{align*}
Using $\Gamma(e,e)=2\mu_c^2$ and $\Gamma(e,h_c)=\mu_ce^{-x_c}x_c(x_c-2)$, \eqref{eq:decomp} yields
\eqref{eq:eta}. The influence function $h_c$ is determined by the complete degree profile. Tracking only the
two lowest-order statistics of the initial state gives the wrong value $\eta_3\approx-0.398$.

\subsection{The sign of the score determines the peeling outcome}\label{sec:mp-global}
Consider the window $[t_{*,M}-\varepsilon_M,t_{*,M}+\varepsilon_M]$ around the bottleneck. For all times
inside the window at which the process has not yet stopped (including the stopping time itself), the local
noise freezing gives
\begin{equation}
Q_n=x_c\sqrt M R_M+M[\ell_M(n/M)-\ell_M(t_{*,M})]+o_p(\sqrt M),\label{eq:freeze}
\end{equation}
with the error uniform over these times. The reason is as follows. Inside the window the added martingale
fluctuation is $O_p(\sqrt{M\varepsilon_M})$, and the accumulated error of the centered drift is
$O_p(\varepsilon_M\vartheta_M)$ times a logarithmic factor, both $o_p(\sqrt M)$. The difference between the
backward weights at the entry and the actual heavy-stub weights satisfies a bound of the same order. The
detailed estimates are in Appendix~\ref{app:technical}.

From this one can deduce success or failure from the sign of the score. Fix $\zeta>0$. If $R_M\le-\zeta$ and
the process has still not stopped by the central time, \eqref{eq:freeze} gives $Q<0$, a contradiction, so
peeling fails. If $R_M\ge\zeta$, the lower bound on the local fluid difference is $-O(1)$, so a first stop
inside the window likewise contradicts \eqref{eq:freeze}. Outside the window, until only few edges remain,
the positive gap of the fluid together with the uniform concentration before stopping rule out a first stop.

Finally one must rule out a small core. Let $d\le AM$. If $e$ edges form a subgraph of minimum degree at
least two, their endpoints occupy at most $v=\lfloor ke/2\rfloor$ vertices, hence
\begin{equation*}
\Pr(\exists\text{such }e\text{ edges})
\le \binom{d}{e}\binom{M}{v}(v/M)^{ke}
\le[C_{k,A}(e/M)^{k/2-1}]^e.
\end{equation*}
For $k\ge3$ one can choose $\delta_0>0$ small enough that the sum over $1\le e\le\delta_0M$ is $o(1)$. If
the process crosses the bottleneck yet still fails in the end, the residual part would violate this bound.
Thus the outcome is determined on both sides of the safe interval $\pm\zeta$; letting $\zeta\downarrow0$ and
using that the Gaussian limit has no atoms gives \eqref{eq:transfer}.

\subsection{Transfer from the configuration model to the Plain model}\label{sec:mp-plain}
In the configuration model, call an edge bad if some endpoint is repeated within that edge, and let
$N^{\rm bad}_M$ be the number of bad edges. The edges are independent, and each edge is bad with probability
$\binom{k}{2}/M+O(M^{-2})$, so
\begin{equation*}
N^{\rm bad}_M\Rightarrow\Pois(\Lambda_k),\qquad
\Lambda_k=\alpha_c\binom{k}{2}.
\end{equation*}
Conditioning on $N^{\rm bad}_M=0$ gives exactly the Plain model. The probability of this event tends to
$e^{-\Lambda_k}>0$, so error events of probability $o(1)$ transfer directly. But for the joint Gaussian law
to be preserved after conditioning, it must also be asymptotically independent of $N^{\rm bad}_M$.

Let $Y_M=(X_{\rm one},R_M)$. Consider planting a fixed number $l$ of special edges. Planting changes the
initial quadratic energy by $O_p(l\log M+l^2)$ and changes the difference-bounded initial influence function
by $O_l(1)$. In ordinary edge-deletion steps the partner pool changes by only finitely many stubs, and the
per-step corrections to the projected drift and quadratic variation are $O_l(M^{-1})$; there are at most $l$
special deletion steps. Hence the Gaussian limit of $Y_M$ at the precritical time is unchanged. The detailed
estimates are in Appendix~\ref{app:technical}.

For any bounded continuous function $f$, exchangeability of the element labels gives
\begin{equation*}
\E[f(Y_M)(N^{\rm bad}_M)_l]\longrightarrow
\Lambda_k^l\E f(Y),
\end{equation*}
where $(b)_l$ is the falling factorial. Expanding the indicator of having no bad edges with
inclusion-exclusion, and bounding the tail of the series by $\|f\|_\infty(dp_M)^l/l!$, gives
\begin{equation*}
\E[f(Y_M)\mathbf1_{\{N^{\rm bad}_M=0\}}]
\longrightarrow e^{-\Lambda_k}\E f(Y).
\end{equation*}
Dividing both sides by the probability of having no bad edges transfers the joint Gaussian law to the Plain
model. Combined with the preceding symmetric-difference error, the Plain result in the one-sided case
follows. Methodologically this step is close to existing conditioning strategies for defects\cite{jl}.

\subsection{Proof of Theorem~\ref{thm:sign} and Corollary~\ref{cor:balanced}}\label{sec:signproof}
For a joint measurement invariant under permutation of the element labels, one can first generate the
unsigned graph $\mathcal G$ and then distribute the signs uniformly according to the given composition.
Define the weighted overlap graph
\begin{equation*}
m_{ij}=|S_i\cap S_j|,\quad
\Omega=\sum_{i<j}m_{ij},\quad \Omega_2=\sum_{i<j}m_{ij}^2,\quad
\ell_i=\sum_{j\ne i}m_{ij},\quad\bar\ell=2\Omega/d.
\end{equation*}
Here $\Omega$ is the total overlap weight. For general $d/M\to\alpha>0$, with $\mu=k\alpha$,
\begin{equation}
\Omega_2/M\xrightarrow p\mu^2/2,\quad
M^{-1}\sum_i(\ell_i-\bar\ell)^2\xrightarrow p\mu^2,\quad
\max_i \ell_i=O_p(\log M).\label{eq:selfavg}
\end{equation}
The first relation comes from the concentration of the overlap mean and $\E(\Omega_2-\Omega)=O(1)$. The
second expands into a sum over overlap edges and adjacent edge pairs, then counts by the pattern of shared
elements to obtain a variance of $O(M)$. The third follows from $\ell_i\le k\max_vD_v$.

Write the sign quadratic form as $Y_s=\sum_{i<j}m_{ij}s_i s_j$. The construction has two steps: first
generate independent signs $t_i$ with $\E t_i=\rho_M$, then flip a uniformly chosen subset of the surplus
sign class so that the total sign count equals the given value exactly. Let $\xi_i=t_i-\rho_M$. When
$|\rho|<1$, the number of signs to flip is $O_p(\sqrt d)$. Projecting to correct the total count gives
\begin{equation*}
Y_s-\rho_M^2\Omega=
\rho_M\sum_i(\ell_i-\bar\ell)\xi_i+
\sum_{i<j}m_{ij}\xi_i\xi_j+o_p(\sqrt M).
\end{equation*}
The dependency graph of the main terms on the right has maximum degree $O(\log M)$, so the $r$-th connected
cumulant divided by $M^{r/2}$ tends to zero ($r\ge3$). The linear and quadratic terms are uncorrelated, and
the total variance is
\begin{equation*}
\rho_M^2(1-\rho_M^2)\sum_i(\ell_i-\bar\ell)^2+(1-\rho_M^2)^2\Omega_2.
\end{equation*}
By \eqref{eq:selfavg}, divided by $M$ this tends to $\mu^2(1-\rho^4)/2$. The detailed estimates are in
Appendix~\ref{app:technical}.

Let $\gamma=(d\rho_M^2-1)/(d-1)$. Combining the exact finite-sample conditional mean with the CLT above gives
\begin{equation}
X=\gamma X_{\rm one}+V,\qquad
\E[e^{iuV}\mid\mathcal G]\xrightarrow p
\exp[-u^2(1-\rho^4)/2].\label{eq:quenched}
\end{equation}
The conditional characteristic function is bounded, so the convergence also holds in $L^1$. This proves that
$V$ is asymptotically independent of any graph-measurable vector that already has a joint limit. The endpoint
case $|\rho|=1$ is handled by $\E V^2=1-\gamma^2\to0$, and there is no need to establish a uniform CLT on
the rare sign class. This completes the proof of Theorem~\ref{thm:sign}. Combining it with the one-sided
joint limit gives the correlation coefficient $\eta_k\rho^2$, completing the chain of proof of the main
theorem.

When $\rho_M\to0$, on the event $\mathcal S$ whose probability is bounded away from zero, $\gamma X_{\rm
one}$ tends to zero; the $L^1$ error of \eqref{eq:quenched} divided by $\Pr(\mathcal S)$ still tends to
zero. The event is invariant under permutation, so the signs can be randomized legitimately. This proves
Corollary~\ref{cor:balanced}.

\subsubsection{Technical estimates}\label{app:technical}
This subsection gives the key estimates in the main chain of proof, with notation continuing from the main
text and the preceding subsections. All constants may depend on the fixed $k$ and on the distance $\delta$ to
the exhaustion of the endpoints, but not on the target time, the truncation degree, or $M$.

\subsubsection{Error of the initial de-Poissonization}
Let $L=kd$ and $w=\lceil\sqrt M\log M\rceil$. On the deterministic interval $[L-w,L+w]$, construct a forward
martingale in the order in which endpoints are added, and use Doob's maximal inequality to control the
martingale difference between any two times. The event $|N-L|\le w$ has probability tending to one, so this
method covers both $N<L$ and $N\ge L$ simultaneously, without a conditional-independence assumption for
deleting endpoints in reverse.

For $e_M$, the accumulated difference between the average increment and $1$ is at most
$O(w^2/M)=O(\log^2 M)$. For $h_c$, write the average increment as $A_h(l)$. Bounded differences give
$A_h(L)-\E A_h(L)=O_p(M^{-1/2})$; the difference between the Binomial and Poisson expectations on a single
bin is $O(M^{-1})$; each added endpoint changes $A_h$ by at most $C/M$. This again gives an accumulated
drift error of $O_p(\log^2M)$. The second moment of the maximum of the martingale corresponding to $h_c$ is
$O(w)$. For $e_M$, stop the martingale when the maximum degree exceeds $\lceil\log M\rceil$; this moment is
$O(w\log^2M)=o(M)$. The degree increases monotonically as endpoints are added, so a union bound on the
binomial tail at the right endpoint alone rules out the bad event of truncation failure.

For the triangular array, first take the drift coefficient as $a_{f,M}=\E[f(P_M+1)-f(P_M)]$ with
$P_M\sim\Pois(\mu_M)$. For $e_M$ this coefficient is exactly one; for $h_c$ it differs from the critical
value $a_f$ by $O(M^{-1/2})$, which after multiplying by $N-L=O_p(\sqrt M)$ contributes only $O_p(1)$.
Applying the ordinary CLT to the Poisson triple $(e_M,h_c,P_M)$ and subtracting $a_f$ times the total-count
coordinate gives all the cross-covariances of \eqref{eq:bridge}. One can also cite directly the existing
conditional CLT, namely Janson's Theorem 2.2\cite{janson2001}. Its conditions all hold here: the Poisson
parameter lies in a positive compact interval, the lattice characteristic function satisfies a uniform
quadratic lower bound, and the required fixed-order moments are bounded. The theorem does not require the
quadratic energy to have a positive exponential moment.

\subsubsection{One-step drift of the backward functional and uniform concentration}
Fix the target step $n$, and let the current step be $s<n$. Each step deletes exactly the $k$ stubs of one
edge, so the number of remaining stubs $N_s=k(d-s)$ is deterministic, and the cancellation in the backward
equation is exact. On the interval under consideration, $N_s\ge k\delta M$ and the maximum degree is at most
$L_M=\lceil\log M\rceil$. For
\[
\phi_l(b)=\Pr\{\operatorname{Bin}(l,b)=j\},\quad j\ge2,
\qquad\text{or}\qquad \phi_l(b)=f_l(b),
\]
we have $\phi_0=\phi_1=0$ and $|\phi_l-\phi_{l-1}|\le2$. Let $b_s=a_M(n/M)/a_M(s/M)$ and
$V_s=\sum_v\phi_{D_v(s)}(b_s)$. Fixing $b_s$, write $\Delta_l=\phi_{l-1}-\phi_l$ and
$A=\sum_vD_v\Delta_{D_v}$; then $|A|\le2N_s$. The forced-deleted leaves contribute nothing to either $V$ or
$A$. Each deleted partner changes $A$ by at most $CL_M$, and the denominator decreases by only one. Exposing
the fixed $k-1$ partners in order gives
\[
\E[\Delta V\mid\mathcal F_s,\ b_s\text{ fixed}]
=\frac{k-1}{N_s}A+O(L_M/M).
\]
This comparison already covers the case where the same vertex is hit several times, and does not need the
hits to be independent.

On an interval at fixed distance from the exhaustion of the endpoints, the first two time derivatives of $b$
are bounded. From the differences of Bernstein polynomials we have $|\partial_b\phi_l|\le CL_M^2$ and
$|\partial_b^2\phi_l|\le CL_M^3$. First expand the time change in the post-deletion state, then switch back
to the pre-deletion state: the second-order remainders over all vertices sum to $O(L_M^3/M)$; at most $k$
vertices change, and the cross terms this produces are $O(L_M^2/M)$. The backward equation cancels the
first-order main terms in both the sampling and the time directions, hence
\[
|\E[V_{s+1}-V_s\mid\mathcal F_s]|\le CL_M^3/M
\le CL_M^4/M.
\]
After $O(M)$ steps the accumulated error is $o(\sqrt M)$. After the process stops, the state and the weights
freeze together. The conditional random difference of each step comes only from at most $k$ deletions, so the
centered martingale difference is at most $4k$ in absolute value.

The weight sum at the starting point satisfies a uniform bounded-difference condition for each independent
endpoint, and its expectation differs from the Poisson thinning value by $O(1)$. The Azuma inequality plus a
union bound over the $O(ML_M)$ targets $(n,j)$ gives the uniform $O_p(\vartheta_M)$ bound in
Section~\ref{sec:mp-stop}. The tail of the Poisson distribution at $j>L_M$, weighted by any fixed polynomial
and multiplied by $M$, still tends to zero. The statement above includes the moment of the first stop.

\subsubsection{Convergence of the dynamic bracket and its separation from the initial state}
Deleting one partner contributes $g_j(b)$ to the projection. When the number of remaining stubs is
$\Theta(M)$, the variance correction from sampling without replacement and from repeated hits is
$O(L_M/M)$; switching the weights of the next moment back to the current weights produces another error of
$O(L_M^2/M)$. So the true one-step conditional variance is
\[
(k-1)\left[\sum_{j\ge2}p_jg_j^2-
 \left(\sum_{j\ge2}p_jg_j\right)^2\right]+O(L_M^2/M).
\]
Replacing the state by the fluid value via uniform concentration gives an error of
$O_p(\vartheta_M L_M^2/M)=o_p(1)$. Truncating the Poisson tail and writing the sum as a Riemann sum gives
\eqref{eq:tau}; the missing final segment of the integral has length $\varepsilon_M\to0$, and the integrand
is bounded, so this part can be ignored.

Let $Z_s$ be the normalized stopped martingale and $C_s=\langle Z\rangle_s$. On the whole sample space
$|\Delta Z_s|\le4k/(x_c\sqrt M)$, and
\[
0\le C_s\le\frac{16k^2d}{x_c^2M},\qquad C_{n_-}\xrightarrow p\tau_k^2.
\]
Applying a one-step conditional Taylor expansion to the compensated exponential $\exp(iuZ_s+u^2C_s/2)$, the
second-order terms cancel and the sum of the third-order remainders is $O(M^{-1/2})$. The modulus of the
exponential is controlled by the deterministic upper bound above. For any $\Psi_M$ measurable with respect to
the initial state and satisfying $|\Psi_M|\le1$, the same telescoping sum gives
\[
\left|\E[\Psi_Me^{iuZ_{n_-}}]
-e^{-u^2\tau_k^2/2}\E \Psi_M\right|
\le C_u\left(M^{-1/2}+\E|C_{n_-}-\tau_k^2|\right)=o(1).
\]
This proves \eqref{eq:stable} uniformly over all initial-state test functions, and does not need sigma-fields
nested in $M$. The independence between the Gaussian limit of the initial state and the dynamic innovation
follows.

\subsubsection{Stopping near the bottleneck and the global outcome}
Write $\mathcal E_n=Q_n-M\ell_M(n/M)$ and $\Delta U_2=U_2-Mu_2$. Sequential partner deletion gives the true
drift
\[
\E[Q_{n+1}-Q_n\mid\mathcal F_n]
=-1+(k-1)\frac{2U_2-Q_n}{Mm_M(n/M)}+O(M^{-1}).
\]
Each deletion changes $2U_2-Q$ by at most three, so this remainder does not require truncating the maximum
degree. After subtracting the fluid drift, the remaining main term is
\[
\frac{k-1}{Mm_M(n/M)}(2\Delta U_2-\mathcal E_n).
\]
Over a window of length $O(M\varepsilon_M)$ it accumulates to $O_p(\varepsilon_M\vartheta_M)=o_p(\sqrt M)$.
Multiplying the increments by $\mathbf1_{\{n<\theta_M\}}$ produces a stopped martingale on the original
probability space, with maximum fluctuation $O_p(\sqrt{M\varepsilon_M})$ inside the window. The last step can
reach $\theta_M$, so this control also covers the stopping moment.

At the entry write $\Delta U_j=U_j(n_-)-Mu_{M,j}(n_-/M)$. The Poisson thinning identity together with
handling the rounding gives
\[
x_c\sqrt M R_M-M\ell_M(t_{*,M})
=-\sum_{j\ge2}f_j(b_-)\Delta U_j+O(1).
\]
Since $|j-f_j(b_-)|\le C\varepsilon_Mj^2$, the above differs from the true entry fluctuation
$-\sum j\Delta U_j$ by $O_p(\varepsilon_M\vartheta_ML_M^3)+o(1)=o_p(\sqrt M)$. This proves
\eqref{eq:freeze}.

The global positivity of the fluid degree-one density comes from the fact that the function
$-\log(1-q)/q^{k-1}$ attains its minimum $\mu_c$ uniquely at $q_c$. Hence the critical fluid has only this
one interior zero before exhaustion. On a compact interval at fixed distance from the exhaustion of the
endpoints, the gap outside the window is at least $cM\varepsilon_M^2-O(\sqrt M)$, so the uniform bound before
stopping rules out a first stop in that interval and outside the window. Inside the window the lowest value
of the fluid difference is $-O(1)$, because the derivative at $t_{*,M}$ is $O(M^{-1/2})$ while the local
second derivative has a positive lower bound.

The rounding in the small-core count can also be handled with uniform constants. Let $v=\lfloor ke/2\rfloor$;
the vertex factor in the binomial upper bound is $e^v(v/M)^{ke-v}$. Its logarithmic derivative
$ke/v-\log(v/M)>0$, so one may replace $v$ by the real number $ke/2$ as an upper bound. Take
$C_{k,A}=e^{1+k/2}A(k/2)^{k/2}$ and choose $\delta_0$ so that $C_{k,A}\delta_0^{k/2-1}<1/2$. Summing the
part with $e\le\log M$ as a geometric series and the part with larger $e\le\delta_0M$ as $2^{-e}$, both
parts are $o(1)$. A path with positive score that reaches the moment with fewer than $\delta_0M$ edges
remaining would violate this bound if it still failed; a path with negative score has its degree-one count
drop to zero before the central time. Finally letting the safe interval $\zeta\downarrow0$ completes the
transfer to the real event.

\subsubsection{Limit after planting finitely many bad edges and the inclusion-exclusion expansion}
At time zero, expose all special edges and the initial degrees of the ordinary edges. Given this information,
the ordinary labeled stubs can be distributed uniformly into the slots of the ordinary edges; as each
ordinary edge is exposed, the assignment of the remaining slots is still uniform. When choosing the leaf with
the smallest vertex label, whether it belongs to a special edge is already predictable. The ordinary partner
pool differs from the full pool by at most $kl$ stubs, so the total variation of the per-step sampling kernel
is $O_l(M^{-1})$, and the corrections to the drift and bracket of the bounded projected increment are of the
same order. There are at most $l$ special steps, whose total contribution is $O_l(1)$.

Thus the initial limit, the uniform concentration, the arrival at the precritical time, and the
compensated-exponential argument all preserve the same limit after planting. Planting preserves the limiting
distribution, not the score of each individual graph; special edges may change the final outcome of an
individual graph.

Let $I_e$ be the indicator that edge $e$ is bad. After fixing $l$ edges with $I_e=1$, the remaining edges are
still independent, so the planting lemma applies. Invariance of the edge labels gives the mixed factorial
moment formula. Since $\sup_M dp_M<\infty$, the absolute value of the $l$-th term of the inclusion-exclusion
series is at most $\|f\|_\infty (dp_M)^l/l!$, and the tail of the series tends to zero uniformly in $M$. This
completes the transfer of the Gaussian law. The symmetric-difference error of the true failure is handled
separately by dividing by the probability of having no bad edges, which tends to a positive number; the two
transfers are done in two steps.

\subsubsection{From iid signs to a fixed total sign count}
Fix a deterministic graph satisfying \eqref{eq:selfavg}, and first handle the case $|\rho|<1$. Then
$\Omega=O(d)$ and $\Omega_\ell=\sum_i\ell_i^2=O(d)$. Take iid signs $t_i$ and let
\[
S=\sum_i t_i,\quad q=|S-d\rho_M|/2,\quad
\epsilon=\operatorname{sign}(S-d\rho_M),\quad J=\{i:t_i=\epsilon\}.
\]
Write $N=|J|$. Draw a uniformly chosen $q$-subset $U$ of $J$ and flip its signs. This mechanism is
equivariant under label permutations, and the output composition is fixed, so the output is exactly the
uniform distribution on that composition. Write $h_i=\sum_jm_{ij}t_j$ and
$I_U=\sum_{i<j,\ i,j\in U}m_{ij}$; expanding edge by edge gives
\[
Y_s-Y_t=-2\epsilon\sum_{i\in U}h_i+4I_U.\label{eq:batchflip}
\]
The edges inside $U$ have both endpoints flipped, so their actual contribution is unchanged; the first term
counts them twice, so the correction $4I_U$ cannot be omitted.

Consider the two cases of $\epsilon$ separately with $T_\epsilon=\sum_i\mathbf1_{\{t_i=\epsilon\}}h_i$, whose
expectation is $p_\epsilon\rho_M\sum_i \ell_i$ with $p_\epsilon=(1+\epsilon\rho_M)/2$. Redrawing $t_z$ changes
$T_\epsilon$ by at most $4\ell_z$, so the Efron--Stein inequality gives $\Var T_\epsilon\le
C\Omega_\ell=O(d)$. Both groups have size $p_\epsilon d+O_p(\sqrt d)$, so for both $\epsilon$ simultaneously
$\bar h_J=\rho_M\bar\ell+O_p(d^{-1/2})$. This also covers the random $\epsilon$ selected by $S$.

Given $t$, sampling without replacement gives the exact variance, together with the expectation bound on the
internal edges:
\begin{align*}
\Var\left(\sum_{i\in U}h_i\given t\right)
&=\frac{q(N-q)}{N(N-1)}\sum_{i\in J}(h_i-\bar h_J)^2
\le\frac{q\Omega_\ell}{N-1},\\
0\le\E[I_U\mid t]&\le\frac{q(q-1)}{N(N-1)}\Omega.
\end{align*}
On the event $q\le K\sqrt d$, $N\ge cd$, apply the conditional Chebyshev and Markov inequalities and then let
$K$ grow; the sampling error is $O_p(d^{1/4})$ and $I_U=O_p(1)$. For a degenerate group the variance is
simply set to zero. Combining $2\epsilon q=S-d\rho_M$ with \eqref{eq:batchflip} gives
\[
Y_s=Y_t-\rho_M\bar\ell(S-d\rho_M)+o_p(\sqrt d),
\]
which is the fixed-total-count projection of Section~\ref{sec:signproof}.

The counting behind the graph self-averaging is as follows. Two distinct overlap pairs sharing only one
element have a constant conditional mean by the symmetry of vertex permutations, so their covariance is zero.
Expanding $\sum_i \ell_i^2$ gives star-triple terms; if two star-triples share two or three elements, their
joint support graph is connected. Exposing independent $k$-subsets in order along a spanning tree, the
product expectations are $O(M^{-3})$ and $O(M^{-2})$ respectively, with $O(d^4)$ and $O(d^3)$ corresponding
terms, so the total variance is $O(M)$. One does not need to assume that all overlaps are mutually
independent.

After projection, the local terms of the linear and quadratic parts number $O(d)$ in total, and the size of
each term as well as the maximum degree of the dependency graph is $O(1+\Delta)$, where
$\Delta=\max_i \ell_i$. The connected cumulant of fixed order $r\ge3$ satisfies
\[
|\kappa_r|\le C_r d(1+\Delta)^{2r-1}=o(M^{r/2}).
\]
The variance converges as computed in Section~\ref{sec:signproof}, and the method of moments gives the normal
limit. For a random graph, take a subsequence of an arbitrary subsequence on which the relevant statistics
converge almost surely and $\Delta\le\log^2d$ holds eventually; arguing graph by graph on this subsequence
yields convergence in probability of the conditional characteristic function. The exact centering coefficient
$\gamma=(d\rho_M^2-1)/(d-1)$ differs from $\rho_M^2$ by a quantity whose product with $\Omega$ is $O(1)$.
The conditional characteristic function is bounded, so the convergence holds in $L^1$, and one may multiply by
any bounded graph-measurable test function to establish joint asymptotic independence. The endpoint
$|\rho|=1$ uses the exact total-variance identity of Section~\ref{sec:signproof}, and the above construction
is not performed on the rare sign class.

\subsubsection{Non-degeneracy of the constants}
The integrand of the dynamic variance is strictly positive on a positive-length interval before the critical
point, so $\tau_k^2>0$. For the initial-state part, the Cauchy--Schwarz inequality gives
$\eta_k^2\le \sigma_0^2/(\sigma_0^2+\tau_k^2)<1$, which guarantees that the density of the selection law is
strictly positive. Since $\beta_k<0$ and failure corresponds to a low score, the threshold is indeed
$z_r=-\beta_k r/\sigma_k$.

All remainders in the finite-dimensional joins are handled by Slutsky's theorem. Conditioning on having no
bad edges preserves $o_p(1)$ because the probability of that event has a positive lower bound. When the score
threshold is replaced by the true failure, the error of the joint CDF is bounded for all $x$ simultaneously
by the probability of the symmetric difference of the events. Finally, dividing by $\Phi(z_r)>0$ gives the
conditional distribution; the limit is continuous, so the CDF convergence is uniform. The convergence of the
conditional mean additionally uses the uniform integrability from $\E X^2=1$ to truncate; the convergence of
the quantile additionally uses the strict positivity of the limiting density. The argument above is for a
fixed-level quantile.
